% ===================================================================
%  TABLEAU N° 7 — Le Village en hiver. — La Maison rustique au
%  printemps.
%  Feuillets 45 a 50 du fac-simile = folios 43 a 48.
%  Correspondance : folio imprime = numero de feuillet - 2.
%
%  Ce tableau est en deux cenes, chacune avec sa propre numerotation
%  d'alineas, repartant de 1 : « Premiere scene » (Le Village en
%  hiver) et « Deuxieme scene » (La Maison rustique au printemps). Le
%  feuillet 45 porte anke la ligne d'apparat « DEUXIEME SERIE » avant
%  le titre du tableau. Le francais est ici beaucoup plus court que
%  l'ido (texto/io/16-tabelo-07.tex, feuillets 49 a 56) : la deuxieme
%  scene s'acheve des l'alinea 16, sans les sous-titres « Farmo-
%  korto. », « La farmisto. » ni « La pultreyo. » que l'ido intercale.
%
%  ATTENTION — LE FAC-SIMILE FRANCAIS EST UNE PRISE DE VUE : le
%  controle 11 devra normaliser chaque page sur sa propre largeur de
%  bloc avant de comparer quoi que ce soit.
% ===================================================================

% --- Feuillet 45 (folio 43, ouverture de tableau) -------------------
\begin{VUpage}[45]{43}
\VUnotes{143.2mm}{%
(1) Voir, pour le detail d'un repas, les tableaux 6 et 13.%
}
%%K t07-apar-1 apar
\VUsaut{4.0mm}
\VUcentre{13.2pt}{90}{DEUXIÈME SÉRIE}
\VUsaut{3.0mm}
\VUfilet{16mm}
\VUsaut{5.6mm}
\VUtitre{15.0pt}{60}{TABLEAU N\textsuperscript{o} 7}
\VUsaut{3.0mm}

%%K t07-tit-1 sub
\VUcentre{12.0pt}{0}{Première scène.}
\VUsaut{3.0mm}

%%K t07-tit-2 sub
\VUtitre{11.4pt}{0}{Le Village en hiver.}
\VUsaut{4.0mm}

%%K t07-c1-01-1 p
\VUblancAlinea
1. --- Pendant les vacances du premier de l'an, j'ai\nl
accompagné mon père à Neuville, petit village où il\nl
avait quelques affaires à traiter.

%%K t07-c1-02-1 p
\VUblancAlinea
2. --- Nous avons pris la \VUgras{diligence}\textsuperscript{(11)} qui fait le\nl
service entre la ville et ce bourg ; mais comme il\nl
faisait très froid, ce voyage dans un véhicule peu\nl
confortable n'a rien eu de bien agréable.

%%K t07-c1-03-1 p
\VUblancAlinea
3. --- Aussi avons-nous éprouvé un réel plaisir\nl
lorsque nous avons vu le \VUgras{conducteur} arrêter sa\nl
voiture sur la place, en face de l'\VUgras{auberge}\textsuperscript{(2)} de\nl
\textit{l'Ours-Blanc}. \VUgras{L'aubergiste}\textsuperscript{(4)} nous attendait sur le\nl
seuil de sa maison en fumant tranquillement sa \VUgras{pipe.}

%%K t07-c1-03-2 p
\VUblancAlinea
Il nous invita à entrer et je ne me fis pas prier\nl
pour m'installer devant le feu de sarment qui pétillait\nl
dans l'âtre. Je me moquais bien alors de l'âpre \VUgras{vent}\nl
\VUgras{du nord} qui faisait grincer la vieille \VUgras{enseigne}\textsuperscript{(3)} et\nl
emportait en noirs tourbillons la \VUgras{fumée}\textsuperscript{(1)} qui sortait\nl
de la cheminée.

%%K t07-c1-04-1 p
\VUblancAlinea
4. --- C'était l'heure du déjeuner. Sans plus atten\cc
dre, nous fîmes honneur au repas simple mais succulent\nl
qui nous fut servi \textsuperscript{(1)}.

%%K t07-c1-05-1 p
\VUblancAlinea
5. --- Après le déjeuner, j'accompagnai mon père,\nl
qui est agent d'assurances, chez ses clients.

%%K t07-c1-05-2 p
\VUblancAlinea
Nous vîmes d'abord le \VUgras{maréchal ferrant}\textsuperscript{(5)}, qui\nl
demeure près de l'auberge. Il était occupé à ferrer
\parplein
\end{VUpage}

% --- Feuillet 46 (folio 44) ------------------------------------------
\begin{VUpage}[46]{44}
%%K t07-c1-05-2 p suite
\VUcontinue
un cheval. J'admirai la force de son compagnon, qui\nl
tenait solidement sur son tablier de cuir le \VUgras{sabot} du\nl
cheval, tandis que le maréchal fixait le \VUgras{fer}\textsuperscript{(6)} à grands\nl
coups de marteau.

%%K t07-c1-06-1 p
\VUblancAlinea
6. --- Nous entrâmes ensuite chez le \VUgras{cordonnier}\textsuperscript{(7)}\nl
du village, le père Jacob. Bien qu'il eût pour enseigne\nl
une superbe \VUgras{botte}, il se contentait de ressemeler\nl
une vieille paire de souliers percés.

%%K t07-c1-06-2 p
\VUblancAlinea
Le brave homme nous dit en riant : « A la ville,\nl
j'étais « cordonnier » et je n'avais pas de clients.\nl
Ici, je suis « savetier » et l'ouvrage abonde. »

%%K t07-c1-07-1 p
\VUblancAlinea
7. --- Il nous parla de son voisin le \VUgras{barbier}\textsuperscript{(8)}, dont\nl
tous les clients étaient mécontents. Ses \VUgras{rasoirs}\nl
étaient toujours ébréchés et ses \VUgras{ciseaux} ne coupaient\nl
jamais. Mais comme il est le seul perruquier du\nl
village, on est bien forcé de se faire raser et couper\nl
les cheveux par lui. C'est d'ailleurs un homme avare\nl
et désagréable.

%%K t07-c1-08-1 p
\VUblancAlinea
8. --- Heureusement, tous les \VUgras{voisins} ne lui ressem\cc
blent pas. Ainsi le \VUgras{charron}\textsuperscript{(9)} est un excellent homme,\nl
travailleur et complaisant. C'est un plaisir de lui\nl
donner une voiture à réparer. Son travail est tou\cc
jours bien fait.

%%K t07-c1-09-1 p
\VUblancAlinea
9. --- Le père Jacob ne se plaignait pas non plus du\nl
\VUgras{sellier}\textsuperscript{(10)}, auquel il aide parfois à coudre ses \VUgras{harnais}\nl
lorsque l'ouvrage presse trop.

%%K t07-c1-09-2 p
\VUblancAlinea
Mon père lui demanda des nouvelles de sa fa\cc
mille. --- « Tout le monde va bien. Ma petite-fille est\nl
à l'\VUgras{école}\textsuperscript{(12)}. Son \VUgras{institutrice}\textsuperscript{(13)} est très contente\nl
d'elle ; c'est une bonne \VUgras{élève}\textsuperscript{(14)}. Ce n'est pas comme\nl
mon mauvais sujet de fils; il ne songe qu'à jouer.\nl
L'\VUgras{instituteur}\textsuperscript{(15)} ne peut rien lui apprendre. »

%%K t07-c1-10-1 p
\VUblancAlinea
10. --- Le brave homme nous raconta ensuite les\nl
nouvelles du village. On allait construire une nouvelle\nl
école, car celle qui était installée à la \VUgras{mairie} deve\cc
nait beaucoup trop petite. C'était d'ailleurs très facile,\nl
car il suffisait d'acheter une partie du jardin du
\parplein
\end{VUpage}

% --- Feuillet 47 (folio 45) ------------------------------------------
\begin{VUpage}[47]{45}
%%K t07-c1-10-1 p suite
\VUcontinue
meunier. Cela ferait bien l'affaire du pauvre homme.\nl
Depuis les froids, les \VUgras{meules} du \VUgras{moulin}\textsuperscript{(16)} ne pou\cc
vaient plus moudre le blé, la \VUgras{roue}\textsuperscript{(17)} étant prise\nl
dans la \VUgras{glace}\textsuperscript{(25)} du \VUgras{ruisseau}\textsuperscript{(23)}.

%%K t07-c1-11-1 p
\VUblancAlinea
11. --- « Mais, à propos de construction, avez-vous\nl
vu notre nouvelle \VUgras{église}\textsuperscript{(18)} ? Elle est très pittoresque\nl
sous son manteau de neige. Les \VUgras{corbeaux}\textsuperscript{(19)}, qui ne\nl
reconnaissent plus leur vieux clocher, se perchent\nl
tristement sur le grand arbre du \VUgras{cimetière}\textsuperscript{(20)}. »

%%K t07-c1-12-1 p
\VUblancAlinea
12. --- Cette idée du cimetière rappela au cordonnier\nl
les gens que mon père avait connus et qui étaient\nl
morts l'année dernière. « Que de \VUgras{tombes}\textsuperscript{(21)}, mon bon\nl
monsieur, sont venues s'ajouter aux autres tombes\nl
sous les \VUgras{cyprès}\textsuperscript{(22)}. La mort a frappé notre vieux\nl
notaire. Tout le village assistait à son enterre\cc
ment. »

%%K t07-c1-13-1 p
\VUblancAlinea
13. --- « Voilà son successeur qui patine sur le ruis\cc
\parplein
\end{VUpage}

% --- Feuillet 48 (folio 46) ------------------------------------------
\begin{VUpage}[48]{46}
%%K t07-c1-13-1 p suite
\VUcontinue
seau. Il est venu tout à l'heure me faire raccommoder\nl
la courroie de son \VUgras{patin}\textsuperscript{(26)} qui était cassée. »

%%K t07-c1-14-1 p
\VUblancAlinea
14. --- « Voilà aussi, sur le bord du ruisseau, le\nl
nouveau \VUgras{garde champêtre}\textsuperscript{(27)}. C'est un ancien gen\cc
darme qui est la terreur des polissons du village. Ils\nl
se sauvent dès qu'ils aperçoivent ses \VUgras{guêtres}\textsuperscript{(29)} et\nl
son \VUgras{képi}\textsuperscript{(28)}. »

%%K t07-c1-15-1 p
\VUblancAlinea
15. --- « Ah ! voici le père Joseph, qui amène une\nl
\VUgras{charrette de fagots}\textsuperscript{(30)} pour le boulanger.

%%K t07-c1-15-2 p
--- « C'est vrai, dit mon père, je reconnais son atte\cc
lage de \VUgras{mules}\textsuperscript{(31)}. J'ai justement besoin de lui parler,\nl
je vais profiter de l'occasion. Au revoir, père Jacob. »

%%K t07-c1-16-1 p
\VUblancAlinea
16. Au moment où nous sortions, les enfants du\nl
village quittaient l'école et s'élançaient en courant\nl
sur la \VUgras{glissoire}\textsuperscript{(33)}, au risque de tomber et de se\nl
casser un membre dans leur \VUgras{chute}\textsuperscript{(34)}. L'un d'eux\nl
prit son \VUgras{traîneau}\textsuperscript{(32)} chez le charron, fit asseoir un\nl
de ses camarades dedans, et partit en courant.

%%K t07-c1-17-1 p
\VUblancAlinea
17. D'autres se précipitèrent vers un tas de\nl
\VUgras{neige}\textsuperscript{(37)} auprès du \VUgras{pont}\textsuperscript{(40)} et se mirent à bom\cc
barder le \VUgras{bonhomme de neige}\textsuperscript{(39)} qu'ils avaient\nl
fait la veille et auquel ils avaient mis un vieux\nl
chapeau, une pipe et un sac d'école en bandoulière.

%%K t07-c1-18-1 p
\VUblancAlinea
18. Ils se souciaient peu du vent glacial et du froid.\nl
La plupart n'avaient même pas de \VUgras{cache-nez}\textsuperscript{(35)} et,\nl
pour se réchauffer après la bataille de \VUgras{boules de}\nl
\VUgras{neige}\textsuperscript{(38)}, il leur suffisait d'une poignée de mar\cc
rons\textsuperscript{(42)} bien chauds achetés à la mère Jacqueline,\nl
la \VUgras{marchande}, qui grelotte sous son vieux \VUgras{châle}\textsuperscript{(43)}\nl
trop mince.

%%K t07-tit-3 sub
\VUsaut{5.4mm}
\VUcentre{12.0pt}{0}{Deuxième scène.}
\VUsaut{3.0mm}

%%K t07-tit-4 sub
\VUtitre{11.4pt}{0}{La Maison rustique au printemps.}
\VUsaut{4.0mm}

%%K t07-c2-01-1 p
\VUblancAlinea
1. --- Malgré tous les plaisirs de l'hiver, c'est avec\nl
une joie profonde que nous saluons le retour du\nl
printemps.

%%K t07-c2-01-2 p
Quel tableau réjouissant que l'aspect d'une cour de\nl
ferme, le soir, au mois de mai !

%%K t07-c2-02-1 p
\VUblancAlinea
2. --- A l'horizon, le \VUgras{soleil}\textsuperscript{(1)} se couche lentement,\nl
inondant la \VUgras{campagne}\textsuperscript{(3)} de ses \VUgras{rayons}\textsuperscript{(2)} de\nl
pourpre. Le \VUgras{laboureur}\textsuperscript{(6)} diligent se hâte de labourer\nl
son \VUgras{champ}\textsuperscript{(4)}.

%%K t07-c2-02-2 p
\VUblancAlinea
Avec sa \VUgras{charrue}\textsuperscript{(5)} attelée d'un vigoureux che\cc
val\textsuperscript{(7)}, il trace le \VUgras{sillon}\textsuperscript{(8)} dans lequel le \VUgras{semeur}\textsuperscript{(9)}\nl
jette à pleine poignée le \VUgras{grain} qui donnera la moisson\nl
blonde.

%%K t07-c2-03-1 p
\VUblancAlinea
3. --- Les \VUgras{faucheurs}\textsuperscript{(11)} armés de leurs \VUgras{faux} ont\nl
coupé (fauché) l'herbe de la prairie, les \VUgras{faneurs} ont\nl
fait sécher le \VUgras{foin}\textsuperscript{(10)} en le retournant avec les\nl
\VUgras{fourches}\textsuperscript{(12)} et les \VUgras{râteaux}, et les voici qui rentrent.

%%K t07-c2-03-2 p
\VUblancAlinea
Les uns marchent devant la lourde voiture traînée\nl
par des bœufs; ils portent leur outil sur l'épaule. Les\nl
autres, plus paresseux, se sont commodément installés\nl
dans le foin parfumé.
\end{VUpage}

% --- Feuillet 49 (folio 47) ------------------------------------------
\begin{VUpage}[49]{47}
%%K t07-c2-04-1 p
\VUblancAlinea
4. --- Ils adressent en passant un salut amical au\nl
\VUgras{jardinier}\textsuperscript{(16)} occupé dans le \VUgras{verger}\textsuperscript{(13)} à tailler les\nl
\VUgras{arbres fruitiers} et à greffer les \VUgras{poiriers}\textsuperscript{(14)}. Les\nl
\VUgras{pruniers}\textsuperscript{(15)} sont en fleurs, et, dans le \VUgras{potager}\textsuperscript{(17)}\nl
bien fumé, les légumes, choux et salades sont déjà\nl
superbes.

%%K t07-c2-04-2 p
\VUblancAlinea
Sur le petit mur, un beau \VUgras{paon}\textsuperscript{(18)} fait la roue\nl
pour faire admirer les superbes plumes de sa queue.

%%K t07-c2-05-1 p
\VUblancAlinea
5. --- Les portes de l'\VUgras{écurie}\textsuperscript{(19)} sont ouvertes et l'on\nl
aperçoit le râtelier et la \VUgras{litière}\textsuperscript{(23)} de la \VUgras{jument}\textsuperscript{(21)}\nl
que le \VUgras{garçon d'écurie}\textsuperscript{(20)} vient de dételer. Le\nl
petit \VUgras{poulain}\textsuperscript{(22)}, si comique sur ses longues jambes,\nl
suit sa mère en gambadant.

%%K t07-c2-06-1 p
\VUblancAlinea
6. --- Près du \VUgras{puits}\textsuperscript{(29)}, les \VUgras{vaches}\textsuperscript{(25)} vont boire\nl
dans l'\VUgras{auge}\textsuperscript{(26)} de bois qui leur sert d'abreuvoir. Le\nl
petit \VUgras{veau}\textsuperscript{(24)} en profite pour téter, et le \VUgras{bœuf}\textsuperscript{(27)},\nl
flairant la bonne odeur du fourrage, mugit joyeuse\cc
ment et dresse fièrement sa tête aux \VUgras{cornes}\textsuperscript{(28)} soli\cc
dement plantées.

%%K t07-c2-07-1 p
\VUblancAlinea
7. --- Tout le monde travaille. La \VUgras{servante}\textsuperscript{(32)} tient\nl
en main la \VUgras{corde}\textsuperscript{(31)} du puits. Elle puise de l'eau dans\nl
un \VUgras{seau}\textsuperscript{(30)} pour donner à boire aux bestiaux. Auprès\nl
de la \VUgras{grange}\textsuperscript{(33)}, un homme ratisse les brins de\nl
paille, et devant la porte de la \VUgras{maison d'habita}\cc
\VUgras{tion}\textsuperscript{(34)}, une vieille \VUgras{fileuse}\textsuperscript{(35)}, avec son \VUgras{rouet} et sa\nl
\VUgras{quenouille}, file le \VUgras{lin} ou le \VUgras{chanvre} comme au\nl
temps jadis.

%%K t07-c2-08-1 p
\VUblancAlinea
8. --- Le \VUgras{fermier}\textsuperscript{(36)} dirige tout ce travail et donne\nl
des ordres à ses serviteurs. Il recommande de rentrer\nl
la \VUgras{herse}\textsuperscript{(40)} et la \VUgras{pioche}\textsuperscript{(39)} que l'on a oubliées près\nl
de la \VUgras{niche}\textsuperscript{(38)} du \VUgras{chien}\textsuperscript{(37)}.

%%K t07-c2-09-1 p
\VUblancAlinea
9. --- Ses cris effarouchent un \VUgras{pigeon}\textsuperscript{(41)} qui fuit à\nl
tire-d'aile vers le \VUgras{pigeonnier}\textsuperscript{(42)}, et les \VUgras{poules}\textsuperscript{(43)},\nl
qui se hâtent de rentrer au \VUgras{poulailler}\textsuperscript{(44)}.

%%K t07-c2-10-1 p
\VUblancAlinea
10. --- Devant la \VUgras{bergerie}\textsuperscript{(48)}, voici le vieux ber\cc
ger\textsuperscript{(45)} qui ramène son \VUgras{troupeau}\textsuperscript{(49)}. Armé de sa
\parplein
\end{VUpage}

% --- Feuillet 50 (folio 48) ------------------------------------------
\begin{VUpage}[50]{48}
%%K t07-c2-10-1 p suite
\VUcontinue
\VUgras{houlette}\textsuperscript{(46)} qui ne le quitte pas plus que sa vieille\nl
\VUgras{blouse}\textsuperscript{(47)} décolorée, il conduit au bercail \VUgras{mou}\cc
\VUgras{tons}\textsuperscript{(50)}, \VUgras{brebis}\textsuperscript{(53)}, \VUgras{agneaux}\textsuperscript{(54)}, \VUgras{chèvres}\textsuperscript{(51)} et\nl
\VUgras{chevreaux}\textsuperscript{(52)}. Son fidèle \VUgras{chien}\textsuperscript{(55)} Labry ferme\nl
la marche et excite par ses aboiements le zèle des\nl
retardataires.

%%K t07-c2-11-1 p
\VUblancAlinea
11. --- Tout ce bruit et tout ce mouvement ne trou\cc
blent pas la placidité de la bonne \VUgras{vache laitière}\textsuperscript{(56)}\nl
qui fait tinter sa \VUgras{clochette}\textsuperscript{(57)} pendant qu'on la trait\nl
devant la porte de l'\VUgras{étable}\textsuperscript{(58)}.

%%K t07-c2-12-1 p
\VUblancAlinea
12. --- Elle a l'air de comprendre qu'il faut qu'elle\nl
donne son lait pour que la \VUgras{fermière}\textsuperscript{(60)} puisse faire\nl
du \VUgras{beurre} dans la \VUgras{baratte}\textsuperscript{(61)} ou confectionner les\nl
excellents \VUgras{fromages}\textsuperscript{(64)} qui s'égouttent sur la\nl
planche posée sur le \VUgras{cuvier}\textsuperscript{(63)}.

%%K t07-c2-13-1 p
\VUblancAlinea
13. --- Toute la \VUgras{laiterie}\textsuperscript{(59)} est tenue très propre\cc
ment par le \VUgras{garçon de ferme}\textsuperscript{(65)}, qui vient de laver\nl
les \VUgras{pots au lait}\textsuperscript{(62)} dans le grand seau qu'il porte à\nl
la main.

%%K t07-c2-14-1 p
\VUblancAlinea
14. --- Avec ses gros sabots, il marche un peu\nl
lourdement, ce qui met en fuite le \VUgras{dindon}\textsuperscript{(68)}, les\nl
\VUgras{canes}\textsuperscript{(66)}, les \VUgras{canards}\textsuperscript{(67)} et les \VUgras{oies}\textsuperscript{(69)}, qui ouvrent\nl
un large \VUgras{bec}\textsuperscript{(70)} et poussent des cris désespérés.

%%K t07-c2-14-2 p
\VUblancAlinea
Le \VUgras{coq}\textsuperscript{(71)}, plus belliqueux, dresse sa \VUgras{crête}\textsuperscript{(72)}\nl
rouge, secoue les plumes de sa \VUgras{queue}\textsuperscript{(73)} et lance un\nl
retentissant « cocorico ».

%%K t07-c2-15-1 p
\VUblancAlinea
15.--- Une \VUgras{poule mère}\textsuperscript{(74)} picore sur le \VUgras{fumier}\textsuperscript{(81)}\nl
avec ses \VUgras{poussins}\textsuperscript{(75)}. Le petit \VUgras{cochon de lait}\textsuperscript{(78)}\nl
se sauve vers sa mère, l'énorme \VUgras{truie}\textsuperscript{(77)} que nous\nl
voyons près de la porcherie.

%%K t07-c2-16-1 p
\VUblancAlinea
16. --- Dans leur niche, les \VUgras{lapins}\textsuperscript{(79)} broutent avi\cc
dement l'\VUgras{herbe}\textsuperscript{(80)} tendre que leur apporte la fille du\nl
fermier. Jeannette aime beaucoup ses lapins et chaque\nl
jour, après avoir été chercher les \VUgras{œufs}\textsuperscript{(76)} frais dans\nl
le poulailler, elle vient rendre visite à ses petits amis.

\VUsaut{6.0mm}
\VUfilet{22mm}
\end{VUpage}
